\begin{solapartat}
\punts{1,25}
\begin{align*}
\alpha\colon\ {}^{238}_{\ 92}\mathrm{U} &\rightarrow {}^{234}_{\ 90}\mathrm{Th} + {}^{4}_{2}\alpha + \gamma\\
\beta^-\colon\ {}^{234}_{\ 90}\mathrm{Th} &\rightarrow {}^{234}_{\ 91}\mathrm{Pa} + {}^{\ \,0}_{-1}e + {}^{0}_{0}\bar{\nu}_e\\
\beta^-\colon\ {}^{234}_{\ 91}\mathrm{Pa} &\rightarrow {}^{234}_{\ 92}\mathrm{U} + {}^{\ \,0}_{-1}e + {}^{0}_{0}\bar{\nu}_e\\
\alpha\colon\ {}^{234}_{\ 92}\mathrm{U} &\rightarrow {}^{230}_{\ 90}\mathrm{Th} + {}^{4}_{2}\alpha + \gamma\\
\alpha\colon\ {}^{230}_{\ 90}\mathrm{Th} &\rightarrow {}^{226}_{\ 88}\mathrm{Ra} + {}^{4}_{2}\alpha + \gamma\\
\alpha\colon\ {}^{226}_{\ 88}\mathrm{Ra} &\rightarrow {}^{222}_{\ 86}\mathrm{Rn} + {}^{4}_{2}\alpha + \gamma
\end{align*}
\destaca{Alternativament} es pot escriure ${}^{\ \,0}_{-1}\beta$ o $e^-$ en lloc de ${}^{\ \,0}_{-1}e$.

\destaca{Alternativament} es pot escriure ${}^{4}_{2}He$ en lloc de ${}^{4}_{2}\alpha$.

No es penalitzarà l'omissió del fotó.

Si no s'indica l'antineutrí cal restar 0,25~p.

No cal que l'estudiant detalli les operacions relacionades amb la conservació de la càrrega i del nombre màssic; són operacions molt senzilles que es poden fer de cap, per tant, només es valorarà si ha estat capaç de deduir correctament els diferents termes de la reacció.
\end{solapartat}

\begin{solapartat}
\punts{0,5} Per trobar l'activitat necessitem el coeficient de desintegració $\lambda$ i el nombre d'àtoms. $\lambda$ l'obtenim del temps de semidesintegració, que és el temps de reduir l'activitat a la meitat: $A(t) = \frac{A_0}{2} = A_0\,e^{-\lambda t_{1/2}}$ i, per tant, $t_{1/2} = \frac{\ln 2}{\lambda}$.

I obtenim:
\begin{align*}
t_{1/2} &= 3,82\times 24\times 3600 = 3,30\times 10^{5}\ \mathrm{s}\\
\lambda &= \frac{\ln 2}{t_{1/2}} = \frac{\ln 2}{3,30\times 10^{5}} = 2,10\times 10^{-6}\ \mathrm{s^{-1}}
\end{align*}
El nombre d'àtoms és:
\[ N = 0,8\ \mathrm{pg} = 8,0\times 10^{-13}\ \mathrm{g}\,\frac{1\ \mathrm{mol}}{222\ \mathrm{g}}\cdot\frac{6,022\times 10^{23}\ \text{àtoms}}{1\ \mathrm{mol}} = 2,17\times 10^{9}\ \text{àtoms} \]
L'activitat és:
\[ A = \lambda N = 2,10\cdot 10^{-6}\times 2,17\times 10^{9} = 4,56\times 10^{3}\ \mathrm{Bq} \]
\punts{0,25} L'activitat per unitat de volum és:
\[ \frac{A}{V} = \frac{4,56\times 10^{3}}{18} = 253\ \mathrm{Bq/m^3} \]
Podem observar que aquest valor és lleugerament inferior al llindar de $300\ \mathrm{Bq/m^3}$ indicat a l'enunciat.

\punts{0,5} L'activitat inicial per unitat de volum és $230\ \mathrm{Bq/m^3}$. L'activitat per unitat de volum al cap d'un temps ve donada per:
\[ A = A_0\,e^{-\lambda t} \rightarrow A/V = A_0/\mathrm{V}\,e^{-\lambda t} \]
i per tant l'activitat al cap de 15 dies és:
\begin{gather*}
t = 15\ \text{dies} = 1,296\times 10^{6}\ \mathrm{s}\\
A/V\,(15\ \text{dies}) = 230\,e^{-\lambda t} = 230\,e^{-2,10\times 10^{-6}\times 1,296\times 10^{6}} = 15,1\ \mathrm{Bq/m^3}
\end{gather*}
\end{solapartat}
